\chapter{Fundamental Limits}
\label{ch:essay-limits}

The Non-Identifiability theorem establishes that certain failures of historical reconstruction are not contingent limitations of particular inference methods but structural properties of the observation map itself. This chapter develops that result in detail, with particular attention to the equivalence-class structure that makes non-identifiability inevitable under specified conditions and to the diagnostic criterion that distinguishes structural limits from correctable observational errors.

An observation map $\obsmap : \History \to \Observable$ transforms complete histories into whatever evidence remains accessible to an observer at a given time. The map is deliberately lossy: it discards information that was present in the history but has not survived to the observable. The question is which histories can be distinguished on the basis of observables alone and which cannot. Two histories $H_1$ and $H_2$ are observationally indistinguishable when $\obsmap(H_1) = \obsmap(H_2)$ despite $H_1 \neq H_2$. Such histories form an equivalence class under the observation map, written $[H]_{\obsmap} = \{ H' : \obsmap(H') = \obsmap(H) \}$, consisting of all histories that produce the same observable outcome as $H$. The size and structure of these equivalence classes determine what can and cannot be recovered through reconstruction.

The elementary direction of the non-identifiability result, Theorem~4.1 of Chapter~4 in \emph{Admissible Reconstruction}, is immediate. If $H_1$ and $H_2$ are distinct members of the same equivalence class, so that $\obsmap(H_1) = \obsmap(H_2)$, then no function $f : \Observable \to \History$ can return both $H_1$ and $H_2$ as the correct reconstruction of the single observable $\obsmap(H_1) = \obsmap(H_2)$, because functions must be single-valued. Any reconstruction procedure faced with this observable must choose one history or the other, or report that the evidence is insufficient to decide, but it cannot return both as the unique correct answer. This establishes impossibility once a collision has been demonstrated, but it does not yet explain why such collisions should be expected to occur or whether they are rare pathologies confined to contrived examples.

The structural direction, Theorem~4.2, addresses this directly and constitutes the deepest result concerning reconstruction limits. Consider observation maps that factor through the reachable-set structure, meaning that the observable is determined by which continuations remain admissible from the history rather than by the full details of the history itself. Formally, $\obsmap(H) = \psi(\reach{H})$ for some function $\psi$, where $\reach{H}$ is the reachable set from the terminal state of history $H$. Such maps record what futures are possible from the current state but not which particular path was taken to arrive there. This class of observation maps is not artificially restrictive; it includes any observation procedure that determines the observable solely by examining what remains admissible continuation rather than by recovering the complete temporal trace of how the system evolved.

Under this condition, whenever two histories $H_1 \neq H_2$ satisfy $\reach{H_1} = \reach{H_2}$, it follows immediately that $\obsmap(H_1) = \psi(\reach{H_1}) = \psi(\reach{H_2}) = \obsmap(H_2)$, producing a collision. The substantive claim is that such collisions are not exceptional. The reachable-set operator $R(\cdot)$ partitions the space of histories into equivalence classes by definition: histories lie in the same class exactly when they share the same reachable set. An admissibility structure induces a trivial partition, in which every equivalence class is a singleton, only when no two distinct histories ever produce the same reachable set. This would require that the reachable set encode the full history rather than merely the current admissible-continuation structure, which contradicts the premise that the observation map factors through reachable sets precisely because it discards historical detail beyond what remains relevant for future admissibility.

The generic case is therefore non-trivial equivalence classes. When histories are individuated by more information than their reachable sets preserve, the pigeonhole principle applies: there are more histories than distinct reachable-set structures, so multiple histories must map to the same structure. These are not edge cases or carefully constructed counterexamples but the expected outcome of any observation map that preserves continuation structure without preserving the full temporal trace of how that structure was established. The structural direction thereby establishes that non-injectivity is a property of the observation map's design rather than a correctable defect.

The equivalence-class perspective clarifies what reconstruction can and cannot achieve. A reconstruction algorithm operating on observables $o = \obsmap(H)$ without access to the hidden history $H$ can, at best, identify the equivalence class $[H]_{\obsmap}$ to which the true history belongs. If that class contains only one element, reconstruction can succeed uniquely. If it contains multiple elements, reconstruction faces an inherent ambiguity: every member of the class is consistent with the observable, and no additional evidence drawn from the same observation map can distinguish them. The question is then whether the ambiguity can be resolved by switching to a less lossy observation map or whether it persists even when the observation preserves all available continuation structure.

This is where the diagnostic corollary, Corollary~4.1 of Chapter~4, becomes essential. Given a collision $\obsmap(H_1) = \obsmap(H_2)$ with $H_1 \neq H_2$, the corollary establishes a criterion for diagnosing the source of reconstruction failure. The diagnosis proceeds by checking whether $\reach{H_1} = \reach{H_2}$. If the reachable sets differ, then the two histories induce different continuation structures, and the observation map has discarded information that a less lossy procedure could in principle have retained. The failure is correctable: a different observation map that preserves more of the reachable-set structure, or that records additional features beyond reachable sets, could distinguish the histories. The reconstruction error is therefore a limitation of the particular observation procedure employed rather than a fundamental property of the histories themselves.

If, conversely, $\reach{H_1} = \reach{H_2}$, then the two histories produce identical continuation structures, and no observation map that factors through reachable-set structure can distinguish them. The equivalence class $[H_1]_{\obsmap}$ contains at least $H_1$ and $H_2$ as distinct elements, and this is not an artifact of the particular $\psi$ chosen but a structural consequence of the histories sharing the same reachable set. Under these conditions, the reconstruction failure is a fundamental limit: the information required to distinguish the histories has ceased to exist in the observable continuation structure, and no refinement of the observation procedure that preserves the factorization through reachable sets can recover it.

The diagnostic corollary thereby provides a principled answer to the question of whether a specific case of reconstruction uncertainty reflects insufficient evidence, inadequate inference methods, or genuine non-recoverability. The distinction is in principle decidable by examining the reachable-set structure directly. In practice, this requires access to the reachable sets $\reach{H_1}$ and $\reach{H_2}$, which may not be directly observable, but the theoretical criterion is clear: check whether the continuation structures coincide. If they do, the limit is structural. If they do not, additional evidence could resolve the ambiguity.

The theorem has broader implications beyond its immediate statement. It establishes that the equivalence classes $[H]_{\obsmap}$ are not arbitrary collections of histories that happen to produce similar outcomes by coincidence but are precisely the sets of histories sharing the same reachable-set structure. The size and internal structure of these classes are therefore determined by the relationship between the space of possible histories and the space of reachable-set structures that those histories can induce. When history space is high-dimensional and reachable-set space is lower-dimensional, equivalence classes are necessarily non-trivial, containing multiple distinct histories. The observation map then performs a many-to-one projection, and reconstruction from observables becomes inherently ambiguous for any history whose equivalence class contains more than one element.

This is the deepest sense in which observables are lossy projections of latent structure. The loss is not merely that some details are omitted, as though a sufficiently careful observer could in principle recover them by examining the evidence more closely. The loss is that entire dimensions of variation in the history space are collapsed into equivalence classes in the observable space, and no observation map preserving only reachable-set structure can recover those collapsed dimensions. The structural non-identifiability established by Theorem~4.2 is therefore not a limitation of observers but a property of the channel through which history is projected into observables. Different histories that induce the same continuation structure have become, from the perspective of any observation map factoring through that structure, the same observable outcome, and the distinction between them is no longer present in the evidence to be discovered.
